\section*{ВВЕДЕНИЕ}
\addcontentsline{toc}{section}{Введение}

\textbf{Актуальность темы исследования.} Сравнительная политология на современном этапе занимает лидирующие позиции в структуре политической науки, выступая не просто автономной дисциплиной, но и ключевой исследовательской стратегией. Институциональные трансформации последних десятилетий, включая завершение «третьей волны» демократизации и нарастание процессов глобализации, требуют постоянного переосмысления методологического инструментария. Развитие теоретических моделей и подходов напрямую влияет на точность анализа современных политических режимов, электоральных систем и публичной политики.

\textbf{Степень научной разработанности.} Теоретико-методологические основания сравнительного анализа закладывались как в трудах классиков политической мысли (политическая морфология Э. Дюркгейма, концепция идеальных типов М. Вебера), так и в работах основоположников «традиционного» и «нового» этапов (Дж. Брайс, Г. Алмонд, С. Роккан). В отечественном дискурсе значительный вклад в систематизацию методологии, теории и методов компаративистики внес Л. В. Сморгунов, чьи исследования связывают теоретический и эмпирический уровни анализа.

\textbf{Объект исследования} --- политические системы, институты и процессы в их сравнительном измерении.

\textbf{Предмет исследования} --- эволюция методологических подходов и специфика сравнительного метода в политической науке.

\textbf{Цель исследования} --- на основе комплексного анализа выявить закономерности исторического развития сравнительной политологии и эксплицировать логику применения сравнительного метода.

Для достижения поставленной цели необходимо решить следующие \textbf{задачи}:
\begin{itemize}
    \item изучить исторические этапы становления сравнительной политологии от традиционного до неоинституционального периода;
    \item рассмотреть сущность сравнения как общего метода познания и его специфику в политических исследованиях;
    \item проанализировать взаимосвязь теории и метода в рамках компаративистского анализа.
\end{itemize}

\textbf{Методологическая база.} Работа опирается на принципы историзма, системный подход, структурно-функциональный анализ, а также на базовые положения макросоциологического и неоинституционального подходов.